% !TeX root = ../../main.tex
\chapter{面向计算精度冗余的低比特分层计算与协同加速方法}\label{chap:acceles}
\section{研究动机}
\subsection{全量高精度稀疏计算的性能开销}
Top-K 稀疏相似度计算需要从大量候选中返回少量最相关条目。常规实现先完成全量 SpMV，再对得分排序或选择。该流程对所有行读取高精度数值，虽然只有少数行最终进入输出，仍付出了完整数据传输与计算成本。将 SpMV 与 Top-K 维护融合能够减少中间结果写回，但若数值表示和候选访问范围不变，仍不能消除全量高精度扫描的主要开销。

在 HBM FPGA 上，数据流执行能降低通用调度开销，而总线有效载荷仍构成约束。高精度数值和坐标占据较大位宽，限制每周期可送入计算核心的非零元素数量；当乘法单元尚有可用资源时，单纯增加计算资源难以改善带宽受限的流水线。因而，需要先降低传输和计算的实际工作量，再匹配对应硬件结构。

\subsection{低比特近似计算的优化机会}
第\ref{chap:analysis}章表明，低比特得分虽然不能无条件复现精确排名，却可用于识别较大的候选集合。基于这一观察，本章区分两个阶段的任务：第一阶段只要求真实 Top-K 尽可能落入候选集合，第二阶段才计算保留候选的高精度得分。将高精度计算从全量行转移到少量候选行，有望在保持较高候选覆盖的同时降低访存开销。

这一分层思想需要解决两个问题。其一，如何选择位宽与剪枝比例，使粗筛的候选遗漏保持在可接受范围内；其二，如何组织低精度连续扫描和高精度随机访问，使第二阶段的少量访问不因寻址和资源重复而失去优势。AccelES 围绕这两点联合设计算法、数据格式和 FPGA 执行架构\cite{jqzhai2025AccelES}。

\section{低比特分层计算方法概述}
AccelES 包括离线数据准备与在线查询。离线阶段对矩阵进行行重排、ReSparse 剪枝、量化与格式转换，将低比特表示和原始 FP32 表示分别存入 HBM。行重排后保留原始行号映射，使输出仍能对应数据库条目。在线阶段将查询量化，多个低精度核心并行完成粗筛，再根据候选行号读取原始 FP32 数据并重算。
\thesisfig[1.00]{acceles/Two_Step_Solve_Topk_SpMV.pdf}{低比特预筛与高精度候选重算的两阶段流程}{fig:acc-workflow}{%
如图\ref{fig:acc-workflow}所示，低比特阶段处理的是完整候选空间的压缩近似表示，高精度阶段处理的是候选集合对应的原始行。第一阶段的 ReSparse 和量化用于减少粗筛开销，第二阶段保留原始数据以避免将剪枝误差直接带入候选得分。系统同时支持只输出近似候选的 AccelES-INT6，以及执行 FP32 重算的 AccelES-FP32 两种配置。前者侧重候选生成吞吐，后者包含候选得分校正，实验分别评估两种配置。
}


设 $\widehat A$ 为剪枝后的矩阵，$Q(\cdot)$ 为量化变换，低比特候选集合为
\begin{equation}
 \mathcal{C}=\operatorname{CandidateSelect}(Q(\widehat A)Q(x),L).
\end{equation}
重算阶段对 $i\in\mathcal{C}$ 求 $s_i=A_ix$，再从候选得分中选择目标结果。其中，CandidateSelect 由各核局部候选的并集实现，候选覆盖同时受近似得分和局部容量约束。

\section{面向 Top-K 计算的低比特近似与剪枝方法}
\subsection{基于低比特量化的近似计算}
AccelES 将 FP32 数值映射为 6 位有符号整数。对待量化对象 $v$，令 $a_{\max}=\max_j|v_j|$，采用
\begin{align}
 S_v&=\frac{2^{B-1}}{a_{\max}},\\
 Q(v_j)&=\clip\left(\round(S_vv_j),-2^{B-1},2^{B-1}-1\right),\quad B=6.
 \label{eq:acc-quant}
\end{align}
全零对象单独映射为零，避免除零。矩阵与查询各自采用对其所有元素一致的正缩放因子，以使理想缩放不改变行间内积排序。舍入和饱和裁剪仍会引入非均匀误差，因此实际候选质量需要测量。

低比特量化带来三方面作用。首先，非零数值字段缩短，提高总线中可容纳的有效记录数。其次，幅值较小的元素在舍入后可能变为零，进一步降低数据量。最后，整数乘法适合 FPGA 的位宽定制与 DSP 复用，减少实现同等并行度所需的资源。但若列索引和行边界未同步压缩，数值字段缩短的收益会受到元数据比例限制，故后续引入专门的数据格式。

实验采用 4 至 6 位量化对比。6 位在所测试候选扩展设置下保持较高候选覆盖，5 位和4 位则可能分别引入最高约10\%和32.5\%的覆盖损失。综合传输成本与候选覆盖，本文选取 INT6 作为实验配置。若数据库向量得分高度接近，或其数值动态范围明显变化，应重新校准位宽与候选预算。

\subsection{基于 ReSparse 的非零元素剪枝}
量化之外，稀疏矩阵中仍可能包含大量对排序贡献较小的非零元素。ReSparse 按行比较元素绝对值，优先保留较大幅值元素，并通过验证查询选择保留比例。图\ref{fig:acc-resparse}给出剪枝示意。这里“小幅值”是启发式的贡献代理，并不意味着该元素对任意查询都不重要；实际影响还取决于查询分量。
\thesisfig[0.62]{acceles/ReSparse.pdf}{ReSparse 按行保留主要非零元素的示意}{fig:acc-resparse}{%
设保留比例为 $\rho$，第 $i$ 行非零数为 $z_i$，保留集合 $\mathcal{S}_i$ 包含该行绝对值最大的若干元素。若所有非零元素都超过用于识别小值的阈值，可保留该行以避免过度剪枝。原方案的阈值基于全矩阵绝对值均值，第四章将进一步讨论其对极高稀疏度数据的局限。保留数量的取整及短行处理应在实现中固定，以使离线编码和验证结果一致。
}


\begin{algorithm}[htbp]
\caption{用于低比特预筛的按行 ReSparse}
\label{alg:resparse}
\begin{algorithmic}[1]
\REQUIRE 稀疏矩阵 $A$，保留比例 $\rho\in[0,1]$，小值判断阈值 $\mu$
\ENSURE 剪枝矩阵 $\widehat A$
\FOR{每个非空行 $A_i$}
\STATE 获取非零索引集合 $\mathcal{N}_i$
\IF{$\min_{j\in\mathcal{N}_i}|A_{ij}|>\mu$}
\STATE $\mathcal{S}_i\leftarrow\mathcal{N}_i$
\ELSE
\STATE 按 $|A_{ij}|$ 降序排列非零元素
\STATE 保留前 $\lceil\rho|\mathcal{N}_i|\rceil$ 个元素，并按数据格式要求处理空行
\ENDIF
\STATE 在 $\mathcal{S}_i$ 上复制数值，其余位置置零
\ENDFOR
\STATE 返回 $\widehat A$
\end{algorithmic}
\end{algorithm}

为分析剪枝误差，记删除集合为 $\mathcal{D}_i$，则
\begin{equation}
 |A_ix-\widehat A_ix|\leq\sum_{j\in\mathcal{D}_i}|A_{ij}||x_j|
 \leq\norm{x}_{\infty}\sum_{j\in\mathcal{D}_i}|A_{ij}|.
 \label{eq:prune-error}
\end{equation}
若所有行的总得分扰动均不超过 $\varepsilon$，且真实排序边界满足 $y_{(K)}-y_{(K+1)}>2\varepsilon$，则该界可给出 Top-K 集合不变的充分条件。该充分条件同时约束排序边界两侧的得分扰动。实际系统并不在线计算这一最坏情况界，而是通过离线候选保持率校准获得可用配置。

原方案使用多个 $K$ 值下的加权候选指标搜索保留比例。实现时可先缓存每行按幅值排序的结果，再对不同 $\rho$ 进行验证，避免重复排序。有限查询上的候选保持率不一定严格随参数单调变化，因此搜索结果仍需通过独立查询验证，以检查所选配置的质量。

\subsection{候选筛选与高精度重算}
系统采用 32 个 Q-core，各自维护局部 Top-16，得到至多512 个候选。将矩阵预先按行随机重排，可减轻高相关结果集中于同一分区的概率，但不能从理论上消除这种情况。若某分区包含超过16 个真实重要候选，则局部截断可能造成遗漏，即使它们的量化得分完全正确。

因此，候选质量由量化、ReSparse 和分区截断共同决定。对每种配置均应计算式\eqref{eq:candidate-recall}，并在最终重算后检查式\eqref{eq:recall}。当 $K$ 增大或候选集中程度增强时，可增加局部保留数量，但会相应增加更新逻辑、候选传输和重算开销。

高精度重算通过原始矩阵恢复候选内的 FP32 得分，使后续排序不再依赖粗筛阶段的量化分数。重算采用相对于 INT6 更高精度的 FP32 运算，仍受浮点舍入误差影响。两阶段系统仍属于近似检索，因为候选遗漏无法由重算修复。

\section{面向分层计算的稀疏数据格式}
\subsection{Ultra-CSR 流式数据格式}
Ultra-CSR 面向第一阶段的顺序扫描，将每个非零记录划分为列索引、量化数值及行切换标志。在列数不超过1024 的配置下，列索引占10 位，INT6 数值占6 位，行标志占1 位，单个记录共17 位。因此512 位传输可容纳30 个记录，剩余2 位未使用。图\ref{fig:acc-ultracsr}对比不同格式的打包方式。
\thesisfig[1.00]{acceles/Ultra-CSR.pdf}{Ultra-CSR 的非零元素打包与位宽利用}{fig:acc-ultracsr}{%
对包内第 $t$ 个非零元素，设行切换标志为 $f_t$，解码器通过前缀和恢复相对行号。若包的起始行号为 $r_0$，在约定标志含义后可写为 $r_t=r_0+\sum_{u\leq t}f_u$。该编码适合按行连续组织的非空行。若连续空行需要保留原始行号，则单个切换位不足以表示跳跃，应采用外部行号映射或额外空行元数据。本文系统面向每行包含有效特征的嵌入数据，预处理仍应检查量化或剪枝后产生的空行。
}


相比为每个元素携带完整行坐标，行标志显著降低重复信息。相比以固定数量行指针填充传输包的设计，按非零元素编码避免了大量无效指针槽位。它使解码与后续乘加都能围绕固定记录数建立并行流水，而无需令每个硬件通道只处理一整行。

\subsection{Random-CSR 随机访问数据格式}
第二阶段依据候选行号访问原始 FP32 数值，数据流由全量顺序扫描转为少量离散行读取。Random-CSR 在 CSR 行指针定位的基础上，将行起始位置、包内偏移和突发长度组织为硬件可直接使用的元数据，从而降低候选行寻址与数据包重组的开销。
\thesisfig[0.73]{acceles/Random-CSR.pdf}{Random-CSR 的行元数据与 FP32 数据包组织}{fig:acc-randomcsr}{%
Random-CSR 包含行元数据和数据包两部分。行元数据记录首包位置及所需包数，数据包保存 FP32 非零值、列索引和边界信息。Reader 根据行号获取包区间，再处理首包中的起始偏移和末包中的结束边界。所使用的数据包支持10 个 FP32 非零元素的并行处理，包格式同时预留控制字段，使数据定位与并行处理保持一致。
}


两种格式共同体现了按阶段选择表示的原则。Ultra-CSR 以扫描吞吐为目标，Random-CSR 以候选定位和读取为目标。保留两份矩阵会增加静态存储及预处理成本，因此该方案更适合数据库相对稳定且存在多次查询的场景；频繁更新数据库时，需额外评估索引和数值副本维护成本。

\section{面向分层计算的 FPGA 异构加速架构}
\subsection{加速器整体架构}
图\ref{fig:acc-arch}展示32 个 Q-core 与1 个 F-core 组成的异构系统。量化矩阵分区存储于32 个 HBM 通道，每个 Q-core 独立处理所属分区。粗筛完成后，各通道 Reader 依据局部候选读取原始数据，仲裁器合并返回的数据流，交给 F-core 重算。完整配置工作频率为200 MHz。
\thesisfig[1.00]{acceles/Architecture_Overview.pdf}{AccelES 的多 Q-core 与 F-core 协同架构}{fig:acc-arch}{%
系统将大量并行资源用于占据主要数据量的低精度阶段，而不是为每个通道配置完整高精度核心。候选数量远小于数据库行数时，单个共享重算核心能够以较低资源代价处理第二阶段。然而，当候选预算增大或候选行显著变长时，F-core 可能转为关键瓶颈，应通过阶段时间和队列占用评估资源比例。
}


\subsection{低精度并行计算核心}
Q-core 包含解码器、元素乘法器、行归约器及局部候选更新单元，如图\ref{fig:acc-qcore}。解码器根据 Ultra-CSR 的行标志恢复坐标。直接顺序累加标志形成较长依赖链，因此采用树形前缀计算降低组合逻辑深度，使一个传输包内的多个非零元素能够并行进入后续流水。
\thesisfig[0.72]{acceles/Q_Core.pdf}{Q-core 的低比特非零元素级流水线}{fig:acc-qcore}{%
每周期30 个非零元素对应30 次查询向量读取。单个 URAM 的双读端口无法直接满足这一需求，系统将查询向量复制15 份，以片上容量换取足够读取端口。查询维度较低且重复访问较多，使这种复制具有可行性。复制的代价随核数和向量维度增长，因而不能直接推广到任意高维检索任务。
}


乘法结果按行归约。由于一行可能跨越多个传输包，归约器保存上一包末行的部分和，并依据新包边界判断是否继续累加。该设计避免以行长度决定单次任务粒度，使多行和长行都能在固定非零元素吞吐下处理。

\subsection{高精度重算核心}
F-core 由 Reader、仲裁、解码、FP32 乘法与归约路径共同组成。32 个 Reader 分别服务对应 HBM 通道，读取候选行后附加通道、包内偏移和结束标记，以保持跨通道合并后的数据归属。轮询仲裁将多条返回流组织为连续输入，同时避免长期忽略某一通道。
\thesisfig[0.72]{acceles/F_Core.pdf}{F-core 的高精度候选重算结构}{fig:acc-fcore}{%
F-core 每周期并行处理10 个 FP32 非零元素，为满足查询向量随机读取，使用5 份片上副本。归约器利用通道标记区分多个未完成行，并通过结束标志判断何时输出完整得分。
}


候选行的读取和包重组由 Reader 完成，如图\ref{fig:acc-reader}所示。该路径先按候选编号定位起始包，再利用包内偏移屏蔽首包中不属于目标行的元素，避免候选定位时的额外细粒度数据拷贝。
\thesisfig[0.72]{acceles/Reader.pdf}{候选行读取、控制标记与数据包重组}{fig:acc-reader}{%
这种组织使高精度路径的计算资源与候选规模相匹配。它同样引入行元数据访问、包边界处理和仲裁开销，这些开销与乘加成本共同构成两阶段计算的在线时间。
}


\subsection{Top-K 排序与结果更新机制}
Q-core 在流式得分产生时维护小规模候选集合。新得分只需与当前候选边界比较，若优于边界则替换相应项，避免将全部得分写回外存后再排序。候选集合较小时，比较和更新逻辑可以用 LUT 实现，并与乘加归约流水重叠。

局部候选维护降低集中排序的规模，但也引入局部截断误差。系统最终汇总32 个局部 Top-16，其512 个条目构成候选并集，并不保证恰为全局 Top-512。提高局部容量可以降低这类误差，却会增大关键路径和资源消耗。因此，局部候选数既是质量参数，也是微架构参数，需要与目标 $K$ 和候选分布共同选择。

\section{实验评估与分析}
\subsection{实验环境与评估方法}
本章采用 AccelES 的32 核配置开展实验评估，结果与已发表研究及配套材料一致\cite{jqzhai2025AccelES}。FPGA 采用 Alveo U280，使用 Vitis HLS 2023.2 实现；CPU 比较使用 Intel Xeon 平台及 MKL、堆选择，GPU 比较使用 H100 PCIe 80 GB 及 cuSPARSE、Thrust。实验同时将量化、ReSparse 和多核近似策略应用于相应 CPU/GPU 配置，以分析算法本身的作用。实验分别比较近似输出与 FP32 重算配置，并通过消融分析区分算法简化和硬件执行的贡献。

\begin{table}[htbp]\centering
\caption{AccelES 使用的真实语料派生稀疏矩阵}\label{tab:acc-data}
\begin{tabular}{lrrrr}\toprule
数据集&行数&列数&非零元素数&密度\\\midrule
Sp.Baidu&640000&900&$8.03\times10^7$&14.02\%\\
Sp.CC\_zh&2000000&900&$3.82\times10^7$&2.12\%\\
Sp.Crawl&2000000&900&$5.56\times10^8$&30.87\%\\
Sp.GloVe&400000&500&$2.77\times10^7$&13.83\%\\
Sp.Wiki News&1000000&900&$2.74\times10^7$&3.04\%\\\bottomrule
\end{tabular}
\end{table}

表\ref{tab:acc-data}列出五个真实语料派生数据集，来源包括 GloVe、中文词向量及 Crawl 等\cite{glove_2014,baidu_encyclopedia_2018,cc_zh_2018,crawl_2018}，稀疏化参考已有方法\cite{sparse_2015}。此外，实验构造六个人工矩阵，覆盖500 万、1000 万和1500 万行，列数为512 至1024，行非零分布采用均匀或 Gamma 分布。全部11 个矩阵的密度范围为1.88\%至30.87\%。

精度实验对每个矩阵使用100 个查询向量，比较 $K=1$ 至100 时真实 Top-K 落入粗筛候选的比例。原文将该指标记作 Precision@K，本章依其集合定义称为候选保持率，以区分大小为512 的预筛集合和最终大小为 $K$ 的输出。离线预处理包括量化、剪枝、格式转换与装载，在线查询评估不重复计入这一成本；Sp.GloVe 的预处理时间为203 s。

\subsection{计算精度与性能评估}
图\ref{fig:acc-precision}展示量化、分区候选维护和 ReSparse 对候选保持率的影响。真实语料实验中，多核近似带来的额外损失较小，加入 ReSparse 后 Precision@100 最高损失约0.87\%；在人工大规模数据上的最大损失约0.95\%。结果表明，所选参数能够在测试负载上保留大部分真实 Top-K，残余质量损失主要体现为粗筛候选遗漏。
\thesisfig[1.00]{acceles/Ablation_Study_on_Precision_and_Density.pdf}{量化、分区候选维护及 ReSparse 的质量与密度消融}{fig:acc-precision}{%
实验结果表明，真实语料平均密度由12.78\%经量化降至11.33\%，进一步经 ReSparse 降至8.22\%。元素删除改善带宽负担，同时使粗筛误差来源增加。两阶段重算将候选内得分计算恢复到 FP32，因而质量分析应同时关注粗筛覆盖和重算后输出，而非仅比较量化值的逐元素误差。
}


图\ref{fig:acc-performance}汇总跨平台性能。采用 FP32 重算的完整配置相对所选 CPU、GPU 和 FPGA 基线分别取得153.3、2.5 和3.4 倍加速。其中，CPU 加速比反映所选处理器、软件实现与实验配置下的性能差异。
\thesisfig[1.00]{acceles/speedup.pdf}{AccelES 跨平台性能与优化配置对比}{fig:acc-performance}{%
非零元素级流水线使 FPGA 性能对行间密度差异相对不敏感，而 CPU 和 GPU 配置在不同数据密度下呈现不同吞吐。对较稠密数据，通用平台可能获得更好的并行利用率；剪枝虽减少元素数，也可能使其利用率下降。因此，计算量减少与运行时间下降之间并不存在固定比例关系。
}


\subsection{关键技术消融分析}
实验结果表明，量化与 ReSparse 应用于 GPU 配置后，平均性能增益约24\%；CPU 多核配置在结合近似策略后同样获得改进。这说明近似算法本身具有跨平台价值，同时也说明完整加速比包含算法和硬件共同作用。对个别数据，进一步稀疏化降低通用处理器效率，反映出执行粒度与负载之间的匹配要求。

\thesisfig[0.62]{acceles/scalability.pdf}{AccelES-INT6 在不同计算核数下的扩展性}{fig:acc-scaling}{%
图\ref{fig:acc-scaling}给出1 至32 核扩展性。随着独立 HBM 通道和计算核增加，系统在所测范围内取得近线性扩展，支持将数据分区与通道绑定的设计。进一步扩展时，跨通道共享资源、候选汇总及高精度路径将成为潜在限制。
}


\thesisfig[0.67]{acceles/memAccess.pdf}{低比特分层计算的数据访问量变化}{fig:acc-memory}{%
图\ref{fig:acc-memory}给出八个矩阵上的访存分析，实验测得平均访问量减少73.5\%。该收益来自低位宽打包、元素剪枝和高精度访问范围缩小的组合，而非单独某一种格式。在真实部署中，还应联合计算原始 FP32 副本、低比特副本及行元数据的静态占用，避免只报告在线流量而忽略容量成本。
}


\subsection{硬件资源与性能开销分析}
\begin{table}[htbp]\centering
\caption{AccelES 资源占用、功耗与能效比较}\label{tab:acc-resource}
\begin{tabular}{lrrr}\toprule
指标&FP32 FPGA 基线&AccelES-INT6&AccelES-FP32\\\midrule
LUT 使用率&46\%&44\%&50\%\\
FF 使用率&39\%&38\%&43\%\\
DSP 使用率&18\%&13\%&14\%\\
BRAM 使用率&28\%&17\%&18\%\\
URAM 使用率&46\%&63\%&65\%\\
频率（MHz）&203&210&200\\
功耗（W）&54&55&58\\
吞吐（GOPS）&33.78&86.43&81.16\\
能效（GOPS/W）&0.63&1.57&1.40\\\bottomrule
\end{tabular}
\end{table}

表\ref{tab:acc-resource}说明，低位宽降低 DSP 需求，但较高非零并行度需要更多查询副本与缓冲，因此 URAM 使用率上升。加入 F-core 后，LUT 和存储资源进一步增加，完整配置频率由 INT6 配置的210 MHz 降至200 MHz。该结果体现分层计算的真实代价：节省算术和传输开销的同时，需要为候选重算和数据路由付出资源。

本系统还有明确适用边界。其列索引和片上查询复制围绕不超过1024 维设计，扩大维度需要调整数据包和缓冲资源；固定局部 Top-16 不适合直接承担任意大的 $K$；数据库更新频繁时，离线格式转换与多副本维护成本不可忽略。此外，即使第一阶段采用低位宽，仍然读取完整候选矩阵，查询成本随数据库规模增长。这一剩余的计算范围冗余构成下一章的直接研究动机。

% !TeX root = ../../main.tex
\section{分层执行的数据一致性与资源约束}
\subsection{离线准备的数据一致性}
两阶段系统同时维护多个矩阵视图，其一致性首先依赖行标识。设原始行号为 $i$，重排后行号为 $\pi(i)$。Q-core 输出的编号通常对应重排后存储位置，Reader 必须按同一编号体系查找 FP32 副本，最终结果再通过逆映射恢复原始条目。若两个副本独立重排，即使粗筛得分正确，重算也可能读取错误对象。

因此，预处理应建立唯一的重排向量，并使量化、ReSparse、格式编码和原始行映射共享这一向量。对每个候选，验证可比较其原始行号、非零数量及列索引序列，而不是仅验证最终输出数量。原始矩阵还应保持只读，使离线剪枝不会意外污染重算数据。

\begin{table}[htbp]\centering\caption{AccelES 各数据视图的一致性要求}
\begin{tabularx}{\textwidth}{p{2.6cm}XX}\toprule
数据视图&应保留的信息&不一致时的影响\\\midrule
原始 FP32 矩阵&原始数值、列号、行标识&重算失去参考语义\\
ReSparse 矩阵&保留位置与原行对应&候选评价与硬件输入不同\\
INT6 矩阵&量化尺度和零元素处理&软硬件得分不一致\\
Ultra-CSR 数据包&行边界和跨包延续&行得分被拆分或错误合并\\
Random-CSR 元数据&首包、偏移、长度&读取越界或混入相邻行\\
逆置换表&原条目与存储行映射&返回 ID 无法对应真实候选\\\bottomrule
\end{tabularx}\end{table}

\subsection{定点累加位宽与溢出控制}
INT6 输入减少乘法器位宽，但累加器仍须容纳整行和。两个有符号6 位整数的乘积绝对值最大可达1024，若一行有 $n_i$ 个参与乘法的元素，则最坏情况绝对和不超过 $1024n_i$。采用 $b_a$ 位有符号累加器时，应满足
\begin{equation}
 2^{b_a-1}-1\geq1024n_i.
\end{equation}
例如仅作上界分析，$n_i\leq1024$ 时使用不少于22 位的有符号累加器可以覆盖这一绝对值上界。该上界给出了累加位宽的正确性约束，具体硬件配置还需结合资源与时序确定。

若实际输入只使用 $[-31,31]$，乘积界可收紧为961，但第三章量化允许 $-32$，应采用与真实编码一致的界。跨包部分和与最终行和必须使用相容位宽，不能只为单包30 个乘积的和分配空间。饱和与回绕也有不同错误语义，应在软硬件仿真中显式区分。

低位宽乘法通过 DSP 打包实现时，需要防止子乘积之间的进位污染。不同打包策略可能对符号扩展和保护位提出额外要求，因此“一个 DSP 支持多路乘法”是具体实现能力，而非任意有符号操作都可无条件复用。资源节省必须与逐元素正确性验证共同检查。

\subsection{跨包边界的状态维护}
连续流式数据包可能包含一行的尾部、若干完整短行以及下一行的头部。归约器应区分完整行输出与未完成行缓存。若当前包首元素继续上一包末行，先合并保存的部分和；若首元素已切换到新行，则先提交上一行，再开始新累加。

设包 $t$ 的末行标识为 $r_t$、未完成和为 $s_t$。包 $t+1$解码后，只有属于同一行的乘积才能与 $s_t$ 合并。一次传输中可能结束多行，候选更新单元的吞吐也需匹配多行完成率。只保证每周期读取30 个元素，却不能及时接收多个行得分时，反压仍会降低实际带宽。

\begin{table}[htbp]\centering\caption{流式边界情况及对应处理}
\begin{tabularx}{\textwidth}{p{3cm}XX}\toprule
情况&处理要求&验证重点\\\midrule
单行跨多个包&持续保存部分和&不能提前输出或丢弃中间包\\
一包内含多个短行&分别归约并提交&多行输出是否堵塞流水\\
包首恰为新行&提交前包末行&避免跨行累加\\
最后一包未填满&屏蔽填充记录&避免产生虚假得分和候选\\
量化后空行&保持行号映射&边界标记能否表达空行跳跃\\
负得分占多数&使用正确候选初值&不可用零值替代负无穷边界\\\bottomrule
\end{tabularx}\end{table}

这些情况可以通过小规模可手算输入验证，独立于大数据集性能测量。小输入应覆盖同一行跨越两个包、整包恰好结束一行以及全负得分等情况。它们的作用是检测控制逻辑，而非证明大规模候选覆盖。

\subsection{局部候选容量的概率解释}
随机重排行缓解重要候选集中在同一核的风险。若作理想化假设，真实 Top-K 均匀独立分配到 $P$ 个分区，则某个分区包含的真实重要候选数近似服从二项分布 $X\sim\mathrm{Binomial}(K,1/P)$。局部容量为 $l$ 时，该分区超容量的概率为
\begin{equation}
 \Pr(X>l)=1-\sum_{j=0}^{l}{K\choose j}(1/P)^j(1-1/P)^{K-j}.
\end{equation}
使用并合界，任意分区超容量概率不超过 $P\Pr(X>l)$。这只是随机分配下的辅助解释，重要候选之间的相关性、实际重排和近似得分排序会使独立假设不成立。

因此，不能用该概率表达式取代实验召回。它可以说明为什么增大局部容量有助于降低分区截断，并帮助选择需要验证的候选数范围。对同一数据库应固定重排种子以便复现，同时使用多个种子检查结果是否对偶然分区敏感。

\subsection{两阶段吞吐与资源分配}
若每个查询的粗筛耗时为 $T_Q$、重算耗时为 $T_F$，不重叠时在线计算成本近似为两者之和。若多个查询能够使用双缓冲形成阶段流水，稳态吞吐可能受 $\max(T_Q,T_F)$限制，但单次查询延迟仍包含其经过各阶段的时间。必须区分本研究已报告的执行方式与这种潜在流水扩展，不能把理想稳态吞吐写成实测结果。

高精度核心共享能够降低资源占用，其合理性依赖候选数据量较小。随着 $L$增加、候选行变长或查询并发提高，F-core 可能饱和。此时新增 Q-core 不会提高系统吞吐，反而扩大候选队列。资源分配宜观察各阶段忙闲比例、队列深度和反压时间，而不是只比较 DSP 利用率。

对宽度为 $W$的传输通道，若候选平均需要 $b$ 个包，每核输出 $l$个候选，则第二阶段至少涉及约 $Plb$个数据包，加上行元数据读取。该数量与第一阶段完整矩阵包数比较，可估计重算是否仍属于小规模任务。真实随机事务延迟还取决于通道并发，不能直接将包数除以峰值带宽得到准确时间。

\subsection{参数校准与泛化验证}
候选预算、保留比例和量化位宽应在校准查询上选择，在独立查询上报告质量。若使用同一批查询调参并汇总，结果可能低估新查询损失。对于数据分布差异较大的语料，应分别记录参数，而不将某个语料上的最优配置当作统一常数。

校准可先固定较宽候选预算验证 INT6，再逐步加入局部候选截断和 ReSparse，最后检查完整组合。顺序加入模块有助于发现损失来源，但每个模块的独立作用还需与全部配置做交叉验证，因为量化与剪枝可能存在相互影响。总体质量损失不能简单等于各模块单独损失之和。

\subsection{更新场景与接口边界}
数据库插入新行时，需要同步生成低比特数据和 FP32 行元数据；修改现有行则可能改变包边界和分区大小。静态构建的连续数据包不一定适合原地插入，因此更可行的扩展包括增量缓冲和周期重构。本文未报告这种动态维护系统的实测性能，属于部署时需要补充的问题。

查询接口应明确返回候选 ID、重算得分还是最终 Top-K。若只传输 ID，主机需要后续访问原始数据时，应计入完整应用成本；若传输得分，则需注明其数值精度和对应缩放。对用户而言，正确的接口规范比单一内核吞吐更能决定系统能否直接替换原有实现。

% !TeX root = ../../main.tex
\section{分层筛选的执行示例与实现边界}
\subsection{候选重算的校正能力与覆盖边界}
为了直观区分分数恢复与候选恢复，表\ref{tab:acc-toy}给出一个人为构造的八行示例。目标为 Top-2，预筛保留3 行；表中分数仅用于解释机制，不来自论文实验或真实量化器测量。真实参考集合为行1 和行2，但行2 的近似分数落到第4 位，未进入预筛集合。

\begin{table}[htbp]\centering\caption{分层筛选的分析性示例：预筛遗漏不能由重算恢复}\label{tab:acc-toy}
\begin{tabular}{cccc}\toprule
行号&FP32 参考分数&近似分数&是否进入 Top-3 预筛\\\midrule
1&0.90&0.89&是\\
2&0.89&0.83&否\\
3&0.88&0.88&是\\
4&0.87&0.87&是\\
5&0.40&0.41&否\\
6&0.30&0.29&否\\
7&0.20&0.21&否\\
8&0.10&0.10&否\\\bottomrule
\end{tabular}\end{table}

对预筛集合 $\{1,3,4\}$执行高精度重算，得到最终结果 $\{1,3\}$，Recall@2 为 $1/2$。结果中每行的分数都可以是正确的 FP32 分数，但结果集合仍不完整。如果把预筛规模提高到4，则该示例可以恢复真实 Top-2。该现象说明，候选容量是质量补偿机制，其有效性取决于分数间隔和扰动大小。

另一个重要情况是近似阶段把行3 排在行1 之前，而二者都在候选中。此时 FP32 重算可以纠正相对顺序。因此，分层执行有意允许候选内部的排序误差，将高精度预算用于最终集合中的次序与分数。它要求保护的是候选覆盖，而不是让低比特阶段复现全部高精度排序。这个目标比逐项精确计算更宽松，也是节省计算的依据。

\subsection{局部候选容量的独立影响}
原架构中每个 Q-core 维护有限局部候选，再由后续模块处理汇集结果。局部选择与全局选择并不等价。假设两个分区各保留1 行，全局目标为 Top-2，分区 A 的分数为0.9 和0.8，分区 B 为0.3 和0.2。即便所有分数均精确，汇集后也只有0.9 和0.3，真实第二名0.8 已被删除。

该反例说明，局部选择的质量依赖分区容量及重要候选的分布。原方案使用32 个 Q-core 和每核16 项，产生最多512 个候选；“最多512 项”描述容量总和，不意味着获得了全局分数最高的512 行。实验中较高覆盖率来自具体数据、划分与容量配置，不能由 $32\times16$ 这一算式推导出保证。

可用每个分区真实 Top-K 数量的分布评估容量压力。设第 $p$ 个分区包含 $k_p$ 个真实目标，局部容量为 $l_p$，若所有分数精确，且次级排序规则一致，则 $l_p\geq k_p$ 对保存该分区目标是充分条件。实际低比特分数还会引入其他行对容量的竞争，所以条件需要进一步考虑近似排序，而非只统计精确目标数量。

为区分两种风险，离线分析可以保存分区编号、真实排名、近似排名和是否入选四类信息。某行若在精确局部 Top-$l_p$中却未进入近似局部集合，主要受数值扰动影响；若在精确局部选择中也被删除，则容量或分区是直接原因。该分解能够指导扩大候选、改变映射或调整剪枝，而不是统一提高所有精度。

\subsection{非零剪枝与行长度变化}
ReSparse 按元素幅值保留部分非零项，其直接目标是减少扫描记录。然而，不同行的原始非零数不同，相同保留比例不一定对应相同绝对长度。例如一行原有4 项，另一行原有40 项，保留一半后分别剩2 项和20 项；长行仍可能跨越多个传输包，短行则受边界信息和对齐成本影响更大。

剪枝后的行还可能变为空行。格式必须保留该行在原矩阵中的身份，否则候选行号与高精度副本不能对应。对空行，数学得分为零，但如果存在负分数，零分行可能影响 Top-K 结果。因此，不能无条件把空行从候选语义中删除。实现是否需要显式处理它们取决于输入约束和输出规则，其适用性由数据条件与输出规则共同决定。

对于含正负值的向量，删除小幅值元素并不保证得分单调下降。若被删项原本贡献负值，近似得分反而增大；若正负项存在抵消，较小的单项变化也可能导致较大的相对误差。因此，幅值剪枝的合理性应通过目标候选覆盖验证，不能把“删除小值”直接解释为“只删除无贡献元素”。

量化顺序也会影响结果。先剪枝再按新最大幅值缩放，与先量化再删除小项可能产生不同的整数分布。第3 章实现应依据原方案的预处理次序复现，同时保留原始矩阵作为 FP32 重算与质量参考。参数校准必须固定这个次序，否则同一保留比例和同一位宽未必对应同一候选集合。

\subsection{流式执行中的依赖与反压}
Q-core 的高并行度只有在数据持续到达、结果能够持续消费时才能转化为吞吐。稀疏行累加依赖行边界信息，Top-K 维护又依赖完整行分数。若排序模块处理速度不足，其输入队列会填满，反压会沿流水返回乘加和读取模块。此时单独增加乘法器数量不能提高完整查询速率。

类似地，多个 Q-core 汇集候选时，汇合点可能形成瞬时突发。共享 F-core 需要根据候选行号读取高精度行，如果某次查询候选行特别长，精排阶段服务时间会增加。合理的流水分析既要看平均候选数量，也要看候选非零总数与请求长度分布。512 个短行与512 个长行可能具有明显不同的重算成本。

可将每条查询的候选重算工作量写为
\begin{equation}
 z_{\mathrm{re}}(q)=\sum_{i\in\mathcal C(q)}\nnz(A_{i,:}),
\end{equation}
其中 $\mathcal C(q)$ 为实际入选候选集合。若使用剪枝后的行长度估计此量，会低估读取原始 FP32 副本的代价；两阶段使用不同矩阵表示，计量时必须明确引用哪份数据。图中报告的总体节省应与这个实际访问集合相联系。

\subsection{吞吐与单查询延迟的区别}
对只有一个请求的场景，查询要顺序经过预筛与精排，完成时间接近两阶段时间之和。对连续请求，若缓冲和数据依赖允许重叠，预筛第 $q+1$ 条查询可以与精排第 $q$ 条查询并行，此时稳态吞吐由较慢阶段决定。共享内存端口、查询缓存复用和结果缓冲容量都可能限制这种重叠。

因此，实验的 QPS 结果需要结合请求提交方式理解。若基线采用批量 GPU 处理，而 FPGA 按流式查询统计，应该确认两者都包含同样的结果生成步骤，并说明批量对延迟的影响。高 QPS 并不自动意味着相同的低单查询延迟，也不意味着所有到达率下的服务尾延迟都会改善。

上述分析刻画了流水执行与查询到达方式的关系，服务级性能仍需在对应负载下验证。后续测量可记录低负载延迟、饱和吞吐和不同到达率下的排队时间，分别判断架构适合交互查询还是批量检索。这样能够把计算能力与服务能力区分开，而不改变本章已完成的硬件验证结论。

\subsection{参数选择与实现可迁移性}
位宽、剪枝比例和候选容量形成联合配置空间。较低位宽可能需要更大候选集合；更强剪枝也可能增加精排量，因而不应分别选择各参数的最优值后直接组合。较稳妥的方法是在固定验证查询上筛除不满足质量目标的配置，再比较这些配置的完整查询时间和资源占用。

在迁移到另一种 FPGA 或不同 HBM 接口时，INT6 配置未必仍最优。数据包位宽、DSP 支持的乘法形状、片上存储粒度和目标频率都会改变资源收益。如果某平台原生支持 INT8 向量指令而不支持 INT6，位级压缩还可能需要解包代价。本文结果支撑的是特定可定制硬件条件下的选择，而非所有处理器上的位宽优越性。

维度增加也影响查询缓存和索引编码。若列索引超过原字段范围，需要扩大索引或引入分块，前述打包密度随之变化。若查询无法完整复制到每个计算单元，访存争用也可能成为新瓶颈。这说明分层思想具有可迁移性，但具体字段、并行度和流水参数仍需要重新设计与验证。

\section{本章小结}
本章针对全量高精度 Top-K 计算的成本，提出由 INT6 粗筛和 FP32 重算组成的分层计算方法，通过 ReSparse 降低粗筛非零规模，并以 Ultra-CSR、Random-CSR 及异构 FPGA 核心匹配两阶段访问模式。已有实验验证了候选保持、性能、访存和硬件资源之间的权衡。该方法将高精度成本集中于少量候选，但粗筛阶段仍遍历整个低比特矩阵。下一章进一步利用数据结构缩减这一计算范围。
